\SourceLeaf{181}sind immer grösser als der angebliche Unterschied zwischen Hypotenuse und einer
Kathete \textls{desjenigen} rechtwinkligen Dreiecks, worin jede der Reihe nach
die andere Kathete ist, demungeachtet erschöpft ihre Summe offenbar die
Hypotenuse $AC$ nie, so gross auch ihre Anzahl genommen wird.

Wir müssten fast bedauern, bei so bekannten und leichten Dingen so lange
verweilt zu haben, wenn nicht diese Schrift, deren Verf. es übrigens wirklich um
Wahrheit zu thun zu sein scheint, durch die Art, wie sie schon vor ihrer
Erscheinung in öffentlichen Blättern angekündigt wurde, eine mehr als gewöhnliche
Aufmerksamkeit auf sich gezogen hätte. Wir bemerken daher hier nur noch, dass
der Verf. nachher auf eine ganz ähnliche, und daher eben so nichtige Art beweisen
will, dass der Winkel $MST$ nicht stumpf sein kann: allein hierbei ist doch ein
wesentlicher Unterschied, weil in der That die Unmöglichkeit dieses Falles in
aller Strenge bewiesen werden kann, welches weiter auszuführen aber hier nicht
der Ort ist.

\Needspace{11\baselineskip}
\begin{center}\rule{3cm}{0.4pt}\end{center}

\begin{center}
\textsc{Bemerkung.}
\end{center}

Diese bereits in Band IV, Seite 364 bis 368 abgedruckte Anzeige ist hier der
Vollständigkeit wegen reproducirt worden. Zwei Figuren, die sich weder beim
Originale noch bei dem ersten Abdrucke finden, sollen das Verständniss des Textes
erleichtern.

\begin{flushright}\textsc{Stäckel.}\end{flushright}

\begin{center}\rule{3cm}{0.4pt}\end{center}
